\section{Results}

\subsection{Widespread intra-urban variation in phenology}
Combinations of demographic rates and interaction coefficients map to distinct regions of coexistence versus priority effects under the Lotka-Volterra framework. The stability-instability inequality structures the boundary between outcomes across the parameter space.

(see Fig. \ref{fig2:variance}).

\begin{figure}[htbp]
  \centering
  \includegraphics[width=0.9\linewidth]{../../temp/fig-main/fig2-variance.png}
  \caption{\textbf{Phenological distributions and variance partitioning of intra-urban and inter-urban components}. Panels (a) and (d) show the proportional contributions of intra-urban (dark gray) and inter-urban (light gray) components to variation in SOS and EOS, respectively, as estimated from mixed-effects models. Each bar associated with the phylogenetic tree represents one species. Panels (b) and (e) illustrate city-level distributions of the SOS (b) and EOS (e) for four representative species which differ in the proportion of variance. The colored horizontal bars indicate the median (dots) and the interquartile range (IQR, bar ends) within each city. The shaded areas represent the median and IQR of city-level median phenological dates. Panels (c) and (f) compare intra-urban (y-axis) and inter-urban (x-axis) IQR for all species, for SOS and EOS respectively. The dashed 1:1 line indicates where intra- and inter-urban variability would be equal.}
  \label{fig2:variance}
\end{figure}


\subsection{Phenology responds to fine-scale temperature heterogeneity}
Species occurrences and niche centroids in climate space show clear environmental differentiation among grassland species along temperature and precipitation gradients, contextualizing how niche separation relates to potential interaction outcomes.

(see Fig. \ref{fig3:nycexample}).

\begin{figure}[htbp]
  \centering
  \includegraphics[width=0.9\linewidth]{../../temp/fig-main/fig3-nycexample.png}
  \caption{\textbf{Phenological responses to urban thermal variation for species and city examples}. Panel (a) and (b) show spatial patterns of mean winter temperature and the SOS for \emph{Acer rubrum} in New York City in 2023. Panel (c) illustrates the marginal relationship between SOS and mean winter temperature for \emph{Acer rubrum} in New York City in 2023. Panels (d) and (e) summarize temperature sensitivities ($\beta_{\mathrm{Temp}}$) estimated from spatial statistical models for all species in New York City, for SOS and EOS, respectively. Species along the x-axis are ordered according to their phylogenetic relationships.}
  \label{fig3:nycexample}
\end{figure}


Species occurrences and niche centroids in climate space show clear environmental differentiation among grassland species along temperature and precipitation gradients, contextualizing how niche separation relates to potential interaction outcomes.

(see Fig. \ref{fig4:heatmap}).

\begin{figure}[htbp]
  \centering
  \includegraphics[width=0.9\linewidth]{../../temp/fig-main/fig4-heatmap.png}
  \caption{\textbf{Phenological responses to intra-urban temperature variation across all species and cities}. Colors indicate the direction of the phenological response, and color intensity reflects the level of statistical significance. Panel (a) and (b) present heatmaps summarizing the phenological response types for each species-city combination for SOS and EOS, respectively. Species along the x-axis are ordered according to their phylogenetic relationships, whereas cities on the y-axis are arranged by latitude.}
  \label{fig4:heatmap}
\end{figure}


\subsection{Phenological responses vary with climatic context}
Species occurrences and niche centroids in climate space show clear environmental differentiation among grassland species along temperature and precipitation gradients, contextualizing how niche separation relates to potential interaction outcomes.

(see Fig. \ref{fig5:marginal}).

\begin{figure}[htbp]
  \centering
  \includegraphics[width=0.9\linewidth]{../../temp/fig-main/fig5-marginal.png}
  \caption{\textbf{Phenological responses to urban thermal variation across CONUS cities and their relationships with cities' climatic contexts}. Panels (a) and (b) map the spatial distribution of the city-level mean temperature sensitivity across all species present in each city, for SOS and EOS, respectively. Colors represent the magnitude of the average phenological temperature effect, and point sizes indicate the variance of species-specific temperature sensitivities within each city. Panels (c-d) show the univariate relationships between phenological temperature sensitivity and the urban climatic context for the 10 most common species across cities. Points and vertical bars indicate species-specific estimates of ($\beta_{\mathrm{Temp}}$ with 95\% credible intervals. Solid lines show fitted ($\beta_{\mathrm{Temp}}$-climatic context relationships obtained by repeatedly resampling posterior draws of $\beta_{\mathrm{Temp}}$ and performing bootstrap resampling; the shaded areas represent the 95\% uncertainty bands across bootstrap replicates.}
  \label{fig5:marginal}
\end{figure}
